% CONFIGURAÇÕES DE APARÃÅ NCIA DO PDF FINAL--------------------------------------
\makeatletter
\hypersetup{%
    portuguese,
    colorlinks=true,   % true: "links" coloridos; false: "links" em caixas de texto
    linkcolor=black,   % Links internos em preto (padrão do modelo UTFPR, sem vínculos coloridos)
    citecolor=black,   % Links de citações/referências em preto
    filecolor=black,   % Links de arquivos em preto
    urlcolor=black,    % URLs em preto
    breaklinks=true,
    pdftitle={\@title},
    pdfauthor={\@author},
    pdfkeywords={RAIS, CAGED, Mercado de TI, Engenharia de Dados, Arquitetura Medalhão, DuckDB, MinIO, Streamlit, Prophet, Forecasting}
}
\makeatother

% ALTERA O ASPECTO DA COR AZUL--------------------------------------------------
\definecolor{blue}{RGB}{41,5,195}

% REDEFINIÇÃÆ’O DE LABELS---------------------------------------------------------
\def\equationautorefname~#1\null{Equa\c c\~ao~(#1)\null}

% CRIA ÍNDICE REMISSIVO---------------------------------------------------------
\makeindex

% "Citado na página X" nas referências (gerado pelo abnt-refinfo do abntex2cite)
% As linhas abaixo estavam suprimindo essa informação e foram desativadas.
% \renewcommand*{\backref}[1]{}
% \renewcommand*{\backrefalt}[4]{}

% CABEÇALHO DAS PÁGINAS---------------------------------------------------------
% "Capítulo N." em caixa normal (a classe já o exibe em itálico) e apenas o
% título do capítulo em maiúsculas, ex.: "Capítulo 1. INTRODUÇÃO".
% Anexado a \textual para sobrescrever o \chaptermark definido pela classe.
\makeatletter
\g@addto@macro\textual{%
    \renewcommand{\chaptermark}[1]{%
        \markboth{\@chapapp\ \thechapter.\ \MakeTextUppercase{#1}}%
                 {\@chapapp\ \thechapter.\ \MakeTextUppercase{#1}}}%
}
\makeatother

% HIFENIZAÇÃÆ’O DE PALAVRAS QUE NÃÆ’O ESTÃÆ’O NO DICIONÁRIO---------------------------
\hyphenation{%
    qua-dros-cha-ve
    Kat-sa-gge-los
}

